\section{Účel systému}
\label{sec:kelvin-purpose}

Hlavním účelem systému Kelvin je podpora výuky, především programovacích předmětů v jazycích C, C++, Python nebo Rust, ale slouží také jako obecný systém pro odevzdávání úloh v dalších předmětech jakou jsou základy počítačové grafiky (ZPG), modelování v grafických aplikacích (MGA) a další.
Systém poskytuje jednotné webové rozhraní, prostřednictvím kterého mohou studenti odevzdávat svá řešení úloh.
Vyučujícím následně umožňuje tato řešení kontrolovat, poskytovat k nim zpětnou vazbu a udělovat bodové hodnocení.
Součástí systému jsou také moduly pro automatické přeložení odevzdaného kódu, jeho spuštění a ověření výstupů.

\subsection{Typický scénář použití}
\label{subsec:kelvin-workflow}

Cyklus jedné úlohy v systému Kelvin zahrnuje několik na sebe navazujících kroků, od vytvoření zadání vyučujícím až po závěrečné manuální hodnocení odevzdaného řešení:

Vyučující nejprve vytvoří úlohu, kde zadá její textový popis, definuje testovací vstupy, očekávané výstupy a nastaví termín odevzdání.
Úloha je pak zpřístupněna přihlášeným studentům okamžitě nebo v předem stanoveném čase.

Student si zadání přečte, během přiděleného času vypracuje řešení ve svém vývojovém prostředí a výsledný zdrojový kód odevzdá prostřednictvím webového rozhraní systému.
Bezprostředně po odevzdání systém spustí automatizované testy.
Systém zkompiluje zdrojový kód a spustí jej proti připraveným testovacím vstupům, následné výsledky jsou porovnány s očekávanými výstupy.
Student okamžitě vidí výsledky testů které prošly i ty, které selhaly a následně může na základě zpětné vazby řešení opravit a odevzdat znovu.

Po uzavření termínu odevzdání přistoupí vyučující k manuálnímu hodnocení.
Prohlíží zdrojový kód odevzdaného řešení přímo v prostředí systému, přidává podle svého hodnocení komentáře ke konkrétním řádkům nebo blokům kódu a na základě celkového posouzení udělí bodové hodnocení.
Student hodnocení a komentáře obdrží a může na ně reagovat, čímž může vznikat diskuse nad zdrojovým kódem a jeho kvalitou.

\subsection{Další formy hodnocení}
\label{subsec:kelvin-other-assessments}

Systém Kelvin je rovněž využíván pro testy, kvízy a jiné formy hodnocení, které nemusí nutně zahrnovat zdrojový kód.
Například pro testování teoretických znalostí nebo algoritmického myšlení prostřednictvím krátkých formulářových nebo textových odpovědí.

\subsection{Klíčové funkce systému}
\label{subsec:kelvin-features}

Systém Kelvin nabízí několik funkcí, které jsou podstatné jak pro běžný provoz, tak pro plánovanou integraci s LLM\@.

\begin{itemize}
    \item \textbf{Automatické evaluace} -- po každém odevzdání systém automaticky zkompiluje studentův kód, spustí jej proti sadě předpřipravených testovacích vstupů a porovná výstupy s očekávanými hodnotami.
    Zdrojový kód je tak testován na jednotném prostředí, což zajišťuje spravedlivé hodnocení funkční správnosti.
    Výsledky jednotlivých testů jsou studentovi zobrazeny okamžitě, což umožňuje rychlou iteraci nad řešením bez nutnosti čekat na zpětnou vazbu od vyučujícího.

    \item \textbf{Automatické udělování bodů} -- na základě výsledků automatických testů může systém přidělit body automaticky podle poměru úspěšných testů.
    Vyučující tak nemusí manuálně hodnotit funkční správnost a může se soustředit na kvalitativní aspekty řešení, jako je čitelnost kódu, vhodnost zvolené datové struktury nebo efektivita algoritmu.

    \item \textbf{Řádkové komentáře} -- vyučující může označit libovolný řádek nebo rozsah řádků zdrojového kódu a připsat k němu připomínku.
    Komentáře jsou studentovi zobrazeny přímo v kontextu kódu, takže okamžitě vidí, ke které konkrétní části řešení se zpětná vazba vztahuje.
    Tato vlastnost je z pohledu integrace LLM zcela zásadní, protože LLM generuje komentáře přirozeně ve formátu \uv{k řádku X patří připomínka Y}.
    Výstup modelu tak lze přímo mapovat na existující datový model systému bez nutnosti zavádět nové typy do databáze.

    \item \textbf{Notifikace} -- systém informuje studenty o nově přidaných komentářích, uděleném hodnocení nebo blížících se termínech odevzdání.
    Vyučující jsou naopak notifikováni o nových odevzdání.
\end{itemize}

\endinput